\documentclass[Total.tex]{subfiles}
\begin{document}
	\chapter{Measure}
		\section{Construction and Extensions of Measures}
			\subsection{Algebras, Semi-Algebras}
				\subsubsection{Algebras}
					\begin{Def}
						Let $X$ be arbitrary set, then a collection $\mathcal{A}$ of subsets of $X$ is an algebra on $X$ if:
						\begin{itemize}
							\item $X,\emptyset\in \mathcal{A}$;(A I)
							\item \begin{gather*}
								A,B\in \mathcal{A}\Rightarrow \begin{cases}
									A\cap B\in \mathcal{A}\qquad (\mbox{A II})\\
									A\cup B\in \mathcal{A}\qquad(\mbox{A III})\\
									A-B\in \mathcal{A}\qquad (\mbox{A IV})
								\end{cases}
							\end{gather*}
						\end{itemize}
					\end{Def}
					
					\begin{proposition}
						For algebra axiom, define
						\begin{gather*}
						 	B\in \mathcal{A}\Rightarrow X-B\in \mathcal{A}\qquad\mbox{(A iv)}
						\end{gather*}
						Then we have
						\begin{gather*}
							\{(AI),(AII),(AIII),(AIV)\}\quad\mbox{equivalent to }\{(AI),(AII),(AII),(Aiv)\}
						\end{gather*}
					\end{proposition}
				\subsubsection{$\sigma$-Algebra}
					\begin{Def}
						An algebra of sets $\mathcal{A}$ is called a \textbf{$\sigma$-algebra} if for any sequence of sets $(A_n)_{n\in \omega}$, we have
						\begin{gather*}
							A_n\in \mathcal{A}\Rightarrow \cup A_n\in \mathcal{A}\qquad \mbox{(SA I)}
						\end{gather*}
					\end{Def}
					Then
					\begin{Def}
						内容...
					\end{Def}
				\subsubsection{Generating}
					\begin{Def}
						Let $X$ be a set. For any family $\mathcal{F}$ of subsets of $X$, there exists a unique $\sigma$-algebra generated by $\mathcal{F}$, denoted by symbol
						\begin{gather*}
							\sigma(\mathcal{F})
						\end{gather*}
					\end{Def}
					\begin{proposition}
						
					\end{proposition}
				\subsubsection{Morphisms}
				\subsubsection{Borel $\sigma$-Algebra}
				\subsubsection{Ring}
				
			\subsection{Measure}
			\subsubsection{Additivity, Subadditivitty;Countable Additivity}
			\begin{Def}
				Let $A$ be a set, take the notion of power set $P(A)$, then, a \textbf{real-valued set function} $\mu$ is defined on a class of sets $\mathcal{A}$ is called \textbf{additive} if
				\begin{gather*}
					内容...
				\end{gather*} 
			\end{Def}
			\subsubsection{Compact Classes, and Countable Additivity}
			\subsubsection{Infinite, and $\sigma$-Infinite Measure}
			\subsection{Lesbegue Measure}
				\subsubsection{Lesbegue Measure}
				\subsubsection{Lesbegue-Stieltjes Meassures}
			\subsection{Monotone and $\sigma$-additive classes of sets}
	\chapter{The Lesbegue Integral}
		\section{Measurable Functions}
	\chapter{Operations on Measures and Functions}
	\chapter{The Space $L^p$ and Spaces of Measures}
		\section{Space $L^p$}
			\begin{Def}
				Let $(X,\mathcal{A},\mu)$ be a measure space, with nonegative measure $\mu$, and $p\in [1,\infty)$. Then denote
				\begin{gather*}
					\mathcal{L}^p(\mu):=\{f:X\rightarrow [-\infty,\infty]||f|^p\mbox{ is }\mu\mbox{-integrable}\}
				\end{gather*}
				
				Define $\mu$-equivalent: The sets $\mathcal{L}^p(\mu)$ are equipped with the equivalent relation, $\mu$-almost everywhere equivalent
				\begin{gather*}
					f\sim g\Leftrightarrow \mu(\{x\in X|f(x)\not=g(x)\})=0
				\end{gather*}
				Then, denote $L^p(\mu)$ as the factor-space of $\mathcal{L}^p(\mu)$ with respect to the equivalence relation. 
				
				The norm structure is specified by
				\begin{gather*}
					\Vert-\Vert_p:
				\end{gather*} 
			\end{Def}
			\subsubsection{Norm $\Vert\cdot\Vert_p$}
			\begin{lemma}
				内容...
			\end{lemma}
			\begin{proof}
				内容...
			\end{proof}
			
			\begin{corollary}
				The function $\Vert\cdot\Vert_p$ is a norm on the space $L^p(\mu)$. 
			\end{corollary}
			\begin{proof}
				内容...
			\end{proof}
			
			\begin{theorem}
				内容...
			\end{theorem}
			\begin{proof}
				内容...
			\end{proof}
			\subsection{Approxiamation of $L^p$.}
			\subsection{The Weak Topology in $L^p$}
		\section{The Hilbert Space $L^2$}
		\section{Uniform Integrability}
		\section{Convergence of Measures}
	\chapter{Integral and Derivative}
		\section{Differentiability of Functions on the Real Line}
	\chapter{Borel,Baire and Souslin Sets}
	\chapter{Measures on Topological Spaces}
		\section{Meaasures on Manifolds}
	\chapter{Weak Convergence of Measure}
	\chapter{Transformations of Measures, and Isomorphisms}
	\chapter{Conditional Measures and Conditional Expectation}
		
\end{document}